\honest{Qualitatively Confused}{Eliezer Yudkowsky}{2008}

``I suggest,'' I begin, that a primary cause of confusing belief, truth and reality is thinking of beliefs as all-or-nothing.

My example is ``the archetypal postmodernist'', who says that ``The Sun goes around the Earth'' is true for Hunga Huntergatherer and ``The Earth goes around the Sun'' is true for Amara Astronomer. \nb{No postmodernist is named or quoted.} No, I reply: different societies have different beliefs. My postmodernist then says that ``I believe `snow is white''' and ``\,`Snow is white' is true'' express the same opinion, and quotes Wittgenstein: ``If there were a verb meaning `to believe falsely', it would not have any significant first person, present indicative.'' \nb{The line is accurate. In the same section Wittgenstein also writes that ``I believe that this is the case'' is used like the assertion ``This is the case'', and yet the hypothesis that I believe it is not used like the hypothesis that it is so. That is close to both sides of this exchange.} The trouble, I say, is that belief and disbelief, being two-valued, look like true and false.

So I switch to probabilities. If I give 70\% to ``snow is white'', a better word than ``true'' for how well that matches reality is ``accuracy''. The natural score is the logarithm of the probability I gave to what actually happened: $-0.51$ bits if snow is white, $-1.73$ if not, and $-0.88$ on average by my own reckoning. My belief that I hold this 70\% belief can be near-certain. So to call the 70\% assignment itself ``true'' is ``a type error''.

If my representations of belief and of belief about belief are different enough, then I am less likely to commit the Mind Projection Fallacy, ``I hope''. Probabilities live in the range from 0 to 1, accuracy from minus infinity to 0, and things themselves are ``just red, or blue, or weighing 20 pounds''. \nb{The post shows that a probabilistic framework keeps these apart. It gives no evidence that people who use it confuse belief and truth less, and its ``I suggest'' and ``I hope'' say as much.}

Uncertainty is a state of mind. A coin is not ``inherently'' 50\% uncertain; it is not even true or false. ``If a coin cannot be true or false, how much less can it assign a 50\% probability to itself?''
